\section{Coupled constraints, parametric supports, and feasible repair}\label{sec:constraints}

The certificates of Sections~\ref{sec:certificates}--\ref{sec:residual}
treat the relaxation construction abstractly. In a solver, each
directional OBBT problem is a linear program whose rows couple original
and auxiliary variables, and these programs are solved again after the
incumbent improves or the box shrinks. This section answers three
questions. What does information stored from earlier solves still prove
after such a change? When can the comparison
inequality~\eqref{eq:matrix-lipschitz} required by
Theorem~\ref{thm:residual} be established although the optimal basis
changes from round to round? How do nonlinear constraints enter a rate
bound?

The tools are linear programming duality, right-hand-side parametric
linear programming~\cite{borrelli2003parametric}, and error bounds for
constraint systems. None of these is new. The results below make precise which conclusions survive a change
of the cutoff, a change of the box with frozen rows, and a change of the
box with rebuilt McCormick rows, and which conclusions need a new
verification. Three kinds of conclusion must be kept apart throughout. A
dual vector bounds a support from above and therefore guarantees a
reduction. A primal witness bounds a support from below and therefore
limits the possible reduction. Only a verified basis region gives the
exact value.

\subsection{Directional support problems}\label{sec:supports-con}

Sections~\ref{sec:supports-con}--\ref{sec:cover-con} consider lifted
relaxations given by linear programs. On a box $B=[\ell,u]$ the relaxation
is a polytope $R(B)\subseteq\R^n\times\R^m$ of lifted points $z=(x,y)$,
where $x$ are the original and $y$ the auxiliary variables, and the
relaxed objective $v$ is affine. Write $d=n+m$. The rows of $R(B)$ include
the box rows $-x_i\le-\ell_i$ and $x_i\le u_i$. For a cutoff $U$, put
$R_U(B)=\{z\in R(B):v(z)\le U\}$. As noted in
Section~\ref{sec:foundations}, compactness of $R(B)$ makes $K_U(B)$ the
projection of $R_U(B)$ onto $x$, so the endpoints of $T_U(B)$ are optimal
values of $2n$ linear programs. Sections~\ref{sec:equality-con}
and~\ref{sec:repair-con} return to the general relaxations of
Section~\ref{sec:foundations}.

Write $p=(-\ell,u)\in\R^{2n}$ and $B(p)$ for the corresponding box. For
$i=1,\ldots,n$ let $c_i=-e_i$, and for $i=n+1,\ldots,2n$ let
$c_i=e_{i-n}$, extended by zeros on the auxiliary coordinates. The exact
Jacobi endpoint map is
\begin{equation}\label{eq:endpoint-map-con}
 F_U(p)_i=\max\{c_i^Tz:z\in R_U(B(p))\},\qquad i=1,\ldots,2n,
\end{equation}
so that $F_U(p)=(-\ell',u')$ when $T_U(B(p))=[\ell',u']$. The box rows give
$F_U(p)\le p$. When the whole construction is valid and monotone, as
assumed in Section~\ref{sec:foundations}, Lemma~\ref{lem:order} gives
$F_U(p)\le F_U(q)$ whenever $p\le q$. All vector inequalities and absolute
values are componentwise. Statements about $F_U$ concern Jacobi rounds;
a sequential round rebuilds the relaxation between directions.

Rows of three kinds react differently when the cutoff or the box changes.
\begin{enumerate}
\item \emph{Fixed rows} have coefficients and right-hand sides that do not
 depend on $B$ or $U$, such as original linear constraints and globally
 valid cuts.
\item \emph{Box rows and the cutoff row} have fixed coefficients and
 right-hand sides that are affine in $(p,U)$.
\item \emph{Rebuilt rows} are recomputed from the current box. For an
 auxiliary variable $y_k$ standing for a product $x_ix_j$ ($i=j$ allowed),
 the McCormick
 rows~\cite{mccormick1976-computability-of-global-solutions-to} are
 \begin{equation}\label{eq:mccormick-rows-con}
 \begin{aligned}
  y_k&\ge\beta x_i+\alpha x_j-\alpha\beta
   &&\text{for }(\alpha,\beta)\in\{(\ell_i,\ell_j),(u_i,u_j)\},\\
  y_k&\le\beta x_i+\alpha x_j-\alpha\beta
   &&\text{for }(\alpha,\beta)\in\{(\ell_i,u_j),(u_i,\ell_j)\}.
 \end{aligned}
 \end{equation}
 Their coefficients on $x$ are box endpoints, their right-hand sides are
 products of endpoints, and the coefficient of $y_k$ is $\pm1$. For $i=j$
 the rows are the two endpoint tangents and the secant of $x_i^2$.
\end{enumerate}
When only the cutoff changes, the rebuilt rows do not change either, and
the support problems form a linear program with a fixed matrix and a
right-hand side that is affine in $U$. When the box changes, rebuilt rows
change their coefficients. Sections~\ref{sec:dual-regions-con}
and~\ref{sec:brackets-con} treat the fixed-matrix case, and
Section~\ref{sec:rebuilt-con} treats rebuilt rows.

Consider therefore the family
\begin{equation}\label{eq:param-lp-con}
 R(\theta)=\{z\in\R^d:Az\le b+E\theta\},\qquad
 h_c(\theta)=\max\{c^Tz:z\in R(\theta)\},\qquad \theta\in\R^k,
\end{equation}
with $A\in\R^{m_A\times d}$, $b\in\R^{m_A}$, $E\in\R^{m_A\times k}$, and
$h_c(\theta)=-\infty$ if $R(\theta)$ is empty. The parameter $\theta$ may
contain the cutoff, box endpoints that enter only through box rows, or
both. For certificates, all data are rational and every inequality is
checked in exact arithmetic. A basis or dual vector returned by a
floating-point solver is a proposal to be checked, not a proof. An exact
certificate concerns the supplied linear program. Validity of its rows
for the original problem and their node-local scope are separate premises.

\subsection{Dual envelopes and exact basis regions}\label{sec:dual-regions-con}

\begin{proposition}[Dual envelope]\label{prop:dual-envelope-con}
Let $\Pi_c=\{\pi\in\R^{m_A}:\pi\ge0,\ A^T\pi=c\}$.
\begin{enumerate}
\item[(a)] For every $\pi\in\Pi_c$ and every $\theta$,
 \begin{equation}\label{eq:dual-envelope-con}
  h_c(\theta)\le\pi^T(b+E\theta).
 \end{equation}
 Consequently, for every finite nonempty $Y\subseteq\Pi_c$,
 $h_c(\theta)\le\min_{\pi\in Y}\pi^T(b+E\theta)$.
\item[(b)] Suppose $\Pi_c\ne\varnothing$, and let
 $\Theta=\{\theta:R(\theta)\ne\varnothing\}$. Then $\Theta$ is a
 polyhedron, and for every $\theta\in\Theta$, $h_c(\theta)$ equals the
 minimum of $\pi^T(b+E\theta)$ over the finitely many vertices $\pi$ of
 $\Pi_c$. In particular, $h_c$ is finite, concave, piecewise affine, and
 continuous on $\Theta$.
\end{enumerate}
\end{proposition}

\begin{proof}
(a) If $R(\theta)$ is empty there is nothing to prove. For
$z\in R(\theta)$, $c^Tz=\pi^TAz\le\pi^T(b+E\theta)$ because $\pi\ge0$.
Maximize over $z$.

(b) $\Theta$ is the projection of the polyhedron
$\{(z,\theta):Az-E\theta\le b\}$ and is therefore a polyhedron. For
$\theta\in\Theta$ the primal problem is feasible and its dual
$\min\{\pi^T(b+E\theta):\pi\in\Pi_c\}$ is feasible, so linear programming
duality gives equal finite optimal values, both attained. The polyhedron
$\Pi_c$ lies in the nonnegative orthant, contains no line, and therefore
has a vertex. A linear function that is bounded below on such a
polyhedron attains its minimum at a vertex, and there are finitely many
vertices. A minimum of finitely many affine functions is concave,
piecewise affine, and continuous.
\end{proof}

Part~(a) uses neither optimality nor primal feasibility nor basis
information at the new parameter. A stored dual vector therefore remains a
valid upper bound however the optimal basis changes. In OBBT, an upper
bound on $h_{c_i}$ that lies below the current endpoint $p_i$ is itself a
valid tightened endpoint and requires no new solve. Storing dual
information from OBBT solves and propagating it after later bound and
cutoff changes is the principle of the Lagrangian variable bounds of
Gleixner et
al.~\cite{gleixner2017-three-enhancements-for-optimization-based}. The
fixed matrix is the premise that makes~\eqref{eq:dual-envelope-con} valid
at every parameter; Section~\ref{sec:rebuilt-con} shows that it fails for
rebuilt rows.

\begin{corollary}[Cutoff multiplier]\label{cor:cutoff-multiplier-con}
Suppose only the cutoff row depends on the parameter: $\theta=U$ and
$E$ is the unit vector of the cutoff row. Let $\pi\in\Pi_c$ be optimal at
$U_0$, and let $\eta\ge0$ be its cutoff multiplier. For every $U$ with
$R(U)\ne\varnothing$,
\begin{equation}\label{eq:cutoff-reduction-con}
 h_c(U)\le h_c(U_0)-\eta(U_0-U).
\end{equation}
If, in addition, $U\le U_0$ and $\pi'\in\Pi_c$ is optimal at $U$ with
cutoff multiplier $\eta'$, then
\[
 \eta(U_0-U)\le h_c(U_0)-h_c(U)\le\eta'(U_0-U).
\]
\end{corollary}

\begin{proof}
Optimality gives $\pi^T(b+EU_0)=h_c(U_0)$, so
$\pi^T(b+EU)=h_c(U_0)-\eta(U_0-U)$, and
Proposition~\ref{prop:dual-envelope-con}(a)
gives~\eqref{eq:cutoff-reduction-con}. Applying the same proposition at
$U_0$ with $\pi'$ gives $h_c(U_0)\le\pi'^T(b+EU_0)=h_c(U)+\eta'(U_0-U)$.
\end{proof}

For a lower cutoff, $\eta(U_0-U)$ is therefore a \emph{guaranteed}
reduction of the support, not an upper bound on the possible reduction.
A zero multiplier guarantees nothing; it does not show that a later
improvement of the incumbent is useless. The right-hand inequality uses a
multiplier at the new cutoff, which is unknown until that problem is
solved; it describes the reduction by the concavity of $h_c$ but does not
predict it. If optimality at $U_0$ has not been verified, the affine
bound remains valid, but subtracting it from a numerical value of
$h_c(U_0)$ does not give a certified reduction.

\begin{proposition}[Exact basis region]\label{prop:basis-region-con}
Let $I$ be a set of $d$ row indices with $A_I$ nonsingular. Put
\[
 z_I(\theta)=A_I^{-1}(b_I+E_I\theta),\qquad
 \pi_I=A_I^{-T}c,
\]
and extend $\pi_I$ by zeros to a vector $\pi\in\R^{m_A}$. If $\pi_I\ge0$,
then on the closed polyhedron
\[
 C_I=\{\theta:Az_I(\theta)\le b+E\theta\}
\]
the point $z_I(\theta)$ is optimal and
\begin{equation}\label{eq:basis-value-con}
 h_c(\theta)=c^Tz_I(\theta)=\pi^T(b+E\theta)\qquad(\theta\in C_I).
\end{equation}
Thus $h_c$ coincides on $C_I$ with an affine function whose gradient is
$E^T\pi$. Conversely, if every $R(\theta)$ with $\theta$ in a set
$\mathcal P$ is a nonempty polytope, then the regions $C_I$ of all bases
$I$ with $A_I^{-T}c\ge0$ cover $\mathcal P$.
\end{proposition}

\begin{proof}
For $\theta\in C_I$ the point $z_I(\theta)$ is feasible by definition.
The extended vector satisfies $\pi\ge0$ and $A^T\pi=A_I^T\pi_I=c$, and the
rows in $I$ are tight, so
$c^Tz_I(\theta)=\pi_I^TA_Iz_I(\theta)=\pi^T(b+E\theta)$.
Proposition~\ref{prop:dual-envelope-con}(a) gives
$h_c(\theta)\le\pi^T(b+E\theta)$, and feasibility gives the reverse
inequality.

For the converse, fix $\theta\in\mathcal P$. A linear objective attains its
maximum over a nonempty polytope at a vertex $z^*$, and the rows active at
$z^*$ span $\R^d$. Optimality gives $c=\sum_{j\in J}\pi_ja_j$ with
$\pi\ge0$ and $J$ a set of active rows, where $a_j^T$ is row $j$ of $A$;
choose such a representation with $|J|$ minimal. The rows in $J$ are
linearly independent: a linear relation among them would allow moving the
coefficients along it until one becomes zero while the others stay
nonnegative. Extend $J$ by active rows to a set $I$ of $d$ linearly
independent active rows and set the added coefficients to zero. Then $A_I$
is nonsingular, $A_I^{-T}c$ is the chosen nonnegative coefficient vector,
$z_I(\theta)=z^*$, and $\theta\in C_I$.
\end{proof}

The first part is the classical description of critical regions in
right-hand-side parametric linear programming, where optimizers and optimal
values are affine on each region~\cite{borrelli2003parametric}. Zero
multipliers, ties between bases, and lower-dimensional regions are
allowed. At a boundary shared by two regions the two affine formulas agree,
but $h_c$ need not be differentiable there; the gradient refers to one
affine piece. Rejecting a singular proposed basis rejects that proposal,
not the relaxation. The polytope premise in the converse is needed. If $s$
is an original variable and $t$ a free auxiliary variable, the set
$\{(s,t):s=0\}$ has attained supports $\max s=\max(-s)=0$, but its two
rows have rank one, so no basis of two rows exists. The polytope premise
holds whenever every original and auxiliary variable has finite bounds, as
in the McCormick relaxations considered here.

A proposed region is verified as follows. If a parameter polytope is the
convex hull of finitely many vertices, it is contained in $C_I$ if and only
if every vertex satisfies the residual inequalities defining $C_I$,
because $C_I$ is convex. For a scalar cutoff, write
$z_I(U)=z^0+Uz^1$ with $z^0=A_I^{-1}b_I$ and $z^1=A_I^{-1}E_I$. Row $j$
becomes
\begin{equation}\label{eq:ranging-con}
 (a_j^Tz^1-E_j)\,U\le b_j-a_j^Tz^0,
\end{equation}
one affine inequality in $U$. Intersecting them gives the exact basis
interval; this is classical right-hand-side ranging. On that interval the
support moves exactly by $\eta$ per unit of cutoff, where $\eta=c^Tz^1$ is
the cutoff multiplier of $\pi$. Outside it the dual expression remains an
upper bound, but the primal formula need not be feasible. All nonbasic
residuals must be checked in the same exact arithmetic as the basis solve.
Exact arithmetic on the basis alone does not certify a point in a model whose
other rows are evaluated in floating point.

\begin{example}[A zero multiplier before a useful tightening]\label{ex:cutoff-con}
Minimize $x_2$ over $0\le x_1\le3$, $0\le x_2\le4$, and
\[
 x_1-x_2\le1,\qquad 4x_1-x_2\le7,\qquad x_1\le 5/2,
\]
with the cutoff $x_2\le U$, $0\le U\le4$. There are no auxiliary
variables. Maximizing $x_1$ gives
\begin{equation}\label{eq:cutoff-example-con}
 h(U)=\min\{1+U,\ 7/4+U/4,\ 5/2\}
 =\begin{cases}
  1+U,&0\le U\le1,\\
  7/4+U/4,&1\le U\le3,\\
  5/2,&3\le U\le4.
 \end{cases}
\end{equation}
On $[0,1]$ the basis consists of $x_1-x_2\le1$ and $x_2\le U$, with primal
point $(1+U,U)$ and multipliers $(1,1)$; the row $4x_1-x_2\le7$ becomes
$4+3U\le7$, so~\eqref{eq:ranging-con} gives exactly $[0,1]$. On $[1,3]$ the
basis consists of $4x_1-x_2\le7$ and $x_2\le U$, with point $(7/4+U/4,U)$
and multipliers $(1/4,1/4)$; the row $x_1-x_2\le1$ requires $U\ge1$ and the
row $x_1\le5/2$ requires $U\le3$. On $[3,4]$ the basis consists of
$x_1\le5/2$ and $x_2\le U$, with point $(5/2,U)$ and multipliers $(1,0)$;
the row $4x_1-x_2\le7$ requires $U\ge3$, and $x_2\le4$ requires $U\le4$.
The three verified intervals cover $[0,4]$, which
proves~\eqref{eq:cutoff-example-con}, and the cutoff slopes are $1$,
$1/4$, and $0$. The three dual vectors give the affine functions $1+U$,
$7/4+U/4$, and $5/2$; their minimum equals $h$ on $[0,4]$, an instance of
Proposition~\ref{prop:dual-envelope-con}(b).

At $U_0=4$ the cutoff multiplier is zero, yet lowering the cutoff to $U=2$
reduces the support from $5/2$ to $9/4$. At $U_0=2$ the old dual remains
valid at $U=1/2$ and gives the bound $7/4+(1/2)/4=15/8$. The extrapolated
primal point $(15/8,1/2)$ is infeasible, since $x_1-x_2=11/8>1$, and the
actual support is $3/2$. The old multiplier guarantees a reduction of
$3/8$, while the actual reduction is $3/4$.
\end{example}

The example is chosen to make these distinctions visible without numerical
ambiguity; it is not a difficult instance.

\subsection{Reconsidering a support after the incumbent improves}\label{sec:brackets-con}

When the incumbent improves from $U_0$ to $U<U_0$ at a fixed box, every
row except the cutoff is unchanged and $R(U)\subseteq R(U_0)$. Stored
primal points and stored dual vectors then bracket each new support.

\begin{proposition}[Primal--dual bracket]\label{prop:bracket-con}
Fix the box and all rows except the cutoff, and write $R(U)$ for the
relaxation with cutoff $U$. Let $c=c_i$ be a signed coordinate objective
with current endpoint $p_i$. Let $W$ be a finite set of complete lifted
points that satisfy every row except possibly the cutoff, and let
$Y\subseteq\Pi_c$ be finite and nonempty. For a new cutoff $U$, put
$W(U)=\{z\in W:v(z)\le U\}$. If $W(U)\ne\varnothing$, then
\begin{equation}\label{eq:bracket-con}
 \max_{z\in W(U)}c^Tz\;\le\;h_c(U)\;\le\;
 \min\Bigl\{p_i,\ \min_{\pi\in Y}\pi^T(b+EU)\Bigr\},
\end{equation}
and consequently
\[
 \max\Bigl\{0,\ p_i-\min_{\pi\in Y}\pi^T(b+EU)\Bigr\}
 \;\le\;p_i-h_c(U)\;\le\;p_i-\max_{z\in W(U)}c^Tz .
\]
\end{proposition}

\begin{proof}
Every point of $W(U)$ is feasible for the new support problem, which
proves the left inequality. The box row gives $h_c(U)\le p_i$, and
Proposition~\ref{prop:dual-envelope-con}(a) gives the dual bound.
Subtracting from $p_i$ gives the second display.
\end{proof}

The left end of the second display is the reduction that can be applied
immediately, without a solve. The right end limits what a new solve of
this one support problem can achieve. If the two ends
of~\eqref{eq:bracket-con} meet, the new support is known exactly. The
bracket concerns one support problem on one frozen relaxation. It says
nothing about the further tightening that rebuilding the relaxation on the
new box may produce. Bounds on all later rounds require a protected box
(Section~\ref{sec:certificates}) or a comparison matrix
(Section~\ref{sec:cover-con}). Bounding the benefit of a whole round
requires information for every direction, not for one favorable direction.

A stored witness that attains the old support certifies that the support
does not change while the cutoff stays at or above its objective value.
The least such value is an exact threshold.

\begin{proposition}[Exact no-change threshold]\label{prop:threshold-con}
Fix the box and all rows except the cutoff, and let $R(U_0)$ be a nonempty
polytope. Let $\mathcal F_0=\{z\in R(U_0):c^Tz=h_c(U_0)\}$ be the optimal
face and
\[
 \underline U=\min\{v(z):z\in\mathcal F_0\}.
\]
For every $U\le U_0$ with $R(U)\ne\varnothing$, $h_c(U)=h_c(U_0)$ if and
only if $U\ge\underline U$.
\end{proposition}

\begin{proof}
If $U\ge\underline U$, a minimizer $z_0$ defining $\underline U$
satisfies the new cutoff, so $h_c(U)\ge c^Tz_0=h_c(U_0)\ge h_c(U)$, the
last inequality because $R(U)\subseteq R(U_0)$. If $U<\underline U$, let
$z$ attain $h_c(U)$, which exists because $R(U)$ is a nonempty polytope.
If $h_c(U)=h_c(U_0)$, then $z\in\mathcal F_0$ and
$v(z)\le U<\underline U$, contradicting the definition of $\underline U$.
Hence $h_c(U)<h_c(U_0)$.
\end{proof}

Computing $\underline U$ requires one additional linear program over the
optimal face, or a lexicographic step after the support solve. Its cost
must be charged like any other solve. A stored optimal witness with a
larger objective value certifies the no-change interval only down to its
own value; exclusion of such a witness does not show that the support
moves, because another optimal point may remain. Exclusion of the
least-objective witness does show this.
Proposition~\ref{prop:threshold-con} is the single-direction,
single-round counterpart of the face thresholds for protected boxes in
Section~\ref{sec:cutoff-thresholds}.

These results sort the directions after an incumbent improvement into
three classes. If $U\ge\underline U$ for direction $i$, the new support
equals the old one; when the current endpoint $p_i$ already equals that
value, for example because the last round applied it, the direction cannot
move in this round. If a stored dual vector gives a bound below $p_i$, that
bound can be applied at once. Otherwise the
bracket~\eqref{eq:bracket-con} is the only information available, and a
budgeted solve is the way to learn more. The classes concern certified
bounds on one support. Whether a solve is worth its cost is a separate
question.

\begin{example}[Example~\ref{ex:cutoff-con} continued]\label{ex:cutoff-bracket-con}
At $U_0=4$ the optimal face is $\{(5/2,x_2):3\le x_2\le4\}$, because a
point with $x_1=5/2$ needs $x_2\ge3$ by the row $4x_1-x_2\le7$. Thus
$\underline U=3$, and Proposition~\ref{prop:threshold-con} shows that the
support is unchanged exactly for $U\in[3,4]$. Lower the cutoff to $U=2$.
The witness $(5/2,3)$ is excluded, so the support decreases. With the
feasible anchor $(0,0)$, the convex combination
$(2/3)(5/2,3)+(1/3)(0,0)=(5/3,2)$ is feasible at the new cutoff, and the
dual vector stored at $U_0=4$ gives $5/2$. The bracket is $[5/3,5/2]$, and
the true value $9/4$ lies inside it. A dual vector that is optimal at a
cutoff in $(1,3)$ is unique and would tighten the upper end to
$7/4+U/4=9/4$.
\end{example}

\subsection{Rebuilt McCormick rows}\label{sec:rebuilt-con}

When the box shrinks and the McCormick rows~\eqref{eq:mccormick-rows-con}
are rebuilt, the constraint matrix itself depends on $p$. Write
\[
 R(p)=\{z:A(p)z\le b(p)\},
\]
where the entries of $A(p)$ are affine in $p$ and those of $b(p)$ are
quadratic in $p$. For products of original variables the coefficients of
the auxiliary variables are constant. This subsection states exactly what
remains valid. Lifted monotonicity, $R(B')\subseteq R(B)$ for
$B'\subseteq B$ with a common objective (Section~\ref{sec:foundations}),
holds for the reference family of McCormick rows and fixed linear rows
(Lemma~\ref{lem:reference-family}).

\begin{lemma}[Frozen rows]\label{lem:frozen-con}
Assume lifted monotonicity. Fix $p_0$, let $A_0z\le b_0$ be the rows of
$R(B(p_0))$, and for $p\le p_0$ define the frozen relaxation
\[
 \widetilde R(p)=\{z:A_0z\le b_0,\ x\in B(p)\}.
\]
\begin{enumerate}
\item[(a)] $R(B(p))\subseteq\widetilde R(p)$. The frozen family has a fixed
 matrix and a right-hand side affine in $p$, so every dual envelope of
 Proposition~\ref{prop:dual-envelope-con} for the frozen family is a valid
 upper bound for the rebuilt supports.
\item[(b)] Let $p_1=F_U(p_0)$, and let $\widetilde F_U$ be the Jacobi map of
 the frozen family at the same cutoff. Then $\widetilde F_U(p)=p_1$ for
 every $p$ with $p_1\le p\le p_0$.
\end{enumerate}
\end{lemma}

\begin{proof}
(a) Lifted monotonicity gives $R(B(p))\subseteq R(B(p_0))$, and the box
rows of $R(B(p))$ give $x\in B(p)$. (b) The frozen cutoff set is
$R_U(B(p_0))\cap\{x\in B(p)\}$. Every point of $R_U(B(p_0))$ has
$x\in T_U(B(p_0))=B(p_1)\subseteq B(p)$, so this set equals
$R_U(B(p_0))$, whose signed supports are $p_1$.
\end{proof}

Part~(a) is why dual information derived from globally valid rows can be
reused at descendant nodes. Part~(b) shows the price: frozen-row OBBT stops
after one round, and its comparison matrix on $[p_1,p_0]$ is zero. Any
progress of iterated OBBT comes from rebuilding. A zero comparison matrix
computed from frozen rows must not be used in~\eqref{eq:matrix-lipschitz}
for the rebuilt map; Theorem~\ref{thm:residual} would then assert that
no motion remains after the first round. Example~\ref{ex:switch-con} below
gives a case where this assertion is false.

Combining old multipliers with the \emph{rebuilt} rows is a different
shortcut, and it is invalid because in general $A(p)^T\pi\ne c$.

\begin{example}[Reusing a dual vector after the box shrinks]\label{ex:rebuilt-dual-con}
Take the square relaxation of Example~\ref{ex:heron} with $r=1$: on
$B=[0,u]$, $u\ge1$, the rows are $s\ge0$, $s\ge2ux-u^2$, and $s\le ux$,
with the cutoff $s\le1$, and the upper support is
$G(u)=(u^2+1)/(2u)$. At $u_0=2$ the optimal basis consists of the tangent
row $4x-s\le4$ and the cutoff, with exact multipliers $\pi=(1/4,1/4)$,
and $G(2)=5/4$. Rebuild on $[0,5/4]$.
\begin{itemize}
\item Old multipliers with rebuilt right-hand sides give
 $\pi^Tb(5/4)=(1/4)(25/16)+1/4=41/64$. This is not a valid bound. The
 point $(x,s)=(1,1)$ satisfies every rebuilt row and the cutoff, and $x=1$
 is an original feasible point with objective $1=U$. The vector $\pi$ is
 not dual feasible for the rebuilt matrix, because
 $2\cdot(5/4)\cdot(1/4)\ne1$.
\item The frozen rows of $[0,2]$ give the valid bound $5/4$ and no
 tightening, as Lemma~\ref{lem:frozen-con}(b) predicts.
\item The corrected bound of Proposition~\ref{prop:corrected-dual-con}
 below gives $71/64$.
\item The exact rebuilt support is $G(5/4)=41/40$.
\end{itemize}
\end{example}

\begin{proposition}[A stored dual vector after a rebuild]\label{prop:corrected-dual-con}
Let $\pi\ge0$, let $p$ be a parameter with $R(p)$ nonempty, and suppose
every $z\in R(p)$ satisfies known bounds $z^L\le z\le z^U$. Put
$q=c-A(p)^T\pi$. Then
\begin{equation}\label{eq:corrected-dual-con}
 h_c(p)\le\pi^Tb(p)+\sum_{j}\max\{q_jz^L_j,\ q_jz^U_j\},
\end{equation}
where terms with $q_j=0$ are zero even if the corresponding bound is
infinite. If $\pi$ is dual feasible for an earlier parameter $p_0$, that
is, $A(p_0)^T\pi=c$, and the change from $A(p_0)$ to $A(p)$ affects only
columns of original variables, then $q=(A(p_0)-A(p))^T\pi$ vanishes on the
auxiliary columns, and only the current box bounds of $x$ enter the sum.
\end{proposition}

\begin{proof}
For $z\in R(p)$, since $\pi\ge0$,
\[
 c^Tz=\pi^TA(p)z+q^Tz\le\pi^Tb(p)+q^Tz
 \le\pi^Tb(p)+\sum_j\max\{q_jz^L_j,\ q_jz^U_j\}.
\]
Maximize over $z$. The last statement follows from $c=A(p_0)^T\pi$.
\end{proof}

In Example~\ref{ex:rebuilt-dual-con}, $q=(1-2\cdot(5/4)/4,\,0)=(3/8,0)$
and $x\le5/4$, so~\eqref{eq:corrected-dual-con} gives
$41/64+(3/8)(5/4)=71/64$. The bound needs no solve and uses the rebuilt
rows. As a function of the upper endpoint $u\in[1,2]$, the corrected
bound with the multipliers from $u_0=2$ is $u-u^2/4+1/4$; for $u>2$ the
residual changes sign and the correction uses the lower bound $0$ of $x$
instead. Applying it repeatedly with
these same multipliers produces valid endpoints $u_{k+1}=u_k-u_k^2/4+1/4$,
which satisfy $u_{k+1}-1=(u_k-1)(3-u_k)/4$ and converge to the exact limit
$1$ at a linear rate of at most $1/2$. Exact OBBT converges to the same
limit quadratically. This is a property of the example; in general the
correction term grows with the change of the coefficients, and the bound is
useful only for moderate changes of the box.

Inequality~\eqref{eq:corrected-dual-con} is the bound of
Proposition~\ref{prop:dual-residual}, written for maximization and applied
to an exact dual vector of a different matrix. For products of original
variables, the auxiliary columns in~\eqref{eq:mccormick-rows-con} have
fixed coefficients $\pm1$, so the correction needs only the current box.
For composite factorable relaxations, products involving auxiliary
variables have box-dependent auxiliary coefficients, and valid bounds on
those auxiliary variables are also needed.

Exact supports with rebuilt rows are again given by basis certificates,
now with parameter-dependent matrices.

\begin{proposition}[Basis certificate with rebuilt rows]\label{prop:rebuilt-basis-con}
Let $A(\cdot)$ and $b(\cdot)$ be continuously differentiable on an open set
$O\subseteq\R^k$, let $I$ be a set of $d$ row indices, and let
$\Pi\subseteq O$ be closed with $A_I(\theta)$ nonsingular for every
$\theta\in\Pi$. Put
\[
 z_I(\theta)=A_I(\theta)^{-1}b_I(\theta),\qquad
 \pi_I(\theta)=A_I(\theta)^{-T}c,
\]
and extend $\pi_I(\theta)$ by zeros to $\pi(\theta)\in\R^{m_A}$. If
$\pi_I(\theta)\ge0$ and $A(\theta)z_I(\theta)\le b(\theta)$ for every
$\theta\in\Pi$, then:
\begin{enumerate}
\item[(a)] $h_c(\theta)=c^Tz_I(\theta)=\pi(\theta)^Tb(\theta)$ for every
 $\theta\in\Pi$;
\item[(b)] $h_I(\theta)=c^Tz_I(\theta)$ is continuously differentiable on
 an open set containing $\Pi$, and
 \begin{equation}\label{eq:lagrangian-derivative-con}
  \partial_jh_I(\theta)=\pi(\theta)^T\bigl(\partial_jb(\theta)
   -\partial_jA(\theta)\,z_I(\theta)\bigr).
 \end{equation}
\end{enumerate}
\end{proposition}

\begin{proof}
(a) At each fixed $\theta\in\Pi$ this is
Proposition~\ref{prop:basis-region-con} for the matrix $A(\theta)$ and the
right-hand side $b(\theta)$. (b) The set where $\det A_I\ne0$ is open and
contains $\Pi$. On it, $A_I^{-1}$ is continuously differentiable by
Cramer's rule, and so are $z_I$ and $h_I$. Differentiating
$A_Iz_I=b_I$ gives
$A_I\,\partial_jz_I=\partial_jb_I-(\partial_jA_I)z_I$, and therefore
$\partial_jh_I=c^TA_I^{-1}(\partial_jb_I-(\partial_jA_I)z_I)
=\pi_I^T(\partial_jb_I-(\partial_jA_I)z_I)$. Rows outside $I$ have zero
multipliers.
\end{proof}

Formula~\eqref{eq:lagrangian-derivative-con} is the derivative of the
Lagrangian $c^Tz+\pi^T(b(\theta)-A(\theta)z)$ with respect to $\theta$ at
the primal--dual pair. For a fixed matrix it reduces to $E^T\pi$, the
derivative of the certified affine value in~\eqref{eq:basis-value-con};
at a boundary shared by two regions the support itself need not be
differentiable. For the McCormick
rows~\eqref{eq:mccormick-rows-con} it is explicit. Let $\sigma=b-Az$
denote the slack of a row. For an underestimator row with corner
$(\alpha,\beta)$,
\[
 \partial\sigma/\partial\alpha=\beta-x_j,\qquad
 \partial\sigma/\partial\beta=\alpha-x_i,
\]
and for an overestimator row the signs are reversed. For a square, both
corner entries are the same endpoint and the two contributions add; the
tangent at an endpoint $e$ gives $\partial\sigma/\partial e=2(e-x)$, and
the secant gives $\partial\sigma/\partial\ell=x-u$ and
$\partial\sigma/\partial u=x-\ell$. Since $\ell_j\le x_j\le u_j$, each
corner entry contributes at most the row's multiplier times the width of
the partner variable, and a box row contributes its multiplier. Because
$p=(-\ell,u)$, derivatives with respect to a lower endpoint change sign.
On a verified region, the matrix of these partial derivatives of the $2n$
supports is their sensitivity matrix. A comparison matrix must bound the
sensitivity matrices of all regions that later rounds can enter.

The maximal region on which a basis is optimal is no longer a polyhedron.
It is described by rational inequalities in $\theta$. A certificate
therefore verifies the conditions of
Proposition~\ref{prop:rebuilt-basis-con} on a chosen closed polyhedron
$\Pi$, such as an interval or a box of parameters. In one parameter this
is exact univariate algebra. In several parameters it requires bounding
rational functions over $\Pi$, for example by exact interval arithmetic on
a subdivision, which is sufficient but can be conservative.

One special case reduces to a fixed matrix. If the coefficient vector of a
rebuilt row changes only by a positive factor, dividing by that factor gives
a row with fixed coefficients and a parameter-dependent right-hand side.
For example, the tangent inequality $2tx-t^2\le0$ obtained from $x^2\le0$
is equivalent to $x\le t/2$ for every tangent point $t>0$. McCormick rows of
products do not have this property in general, because the ratio between
their $x$- and $y$-coefficients changes with the box.

\subsection{A comparison matrix across active-set changes}\label{sec:cover-con}

Theorem~\ref{thm:residual} needs the comparison
inequality~\eqref{eq:matrix-lipschitz} on an invariant region. A
derivative computed at the current basis does not supply it, because later
rounds can enter regions with larger sensitivities. The next theorem
derives a two-sided form of~\eqref{eq:matrix-lipschitz} from a finite cover
by regions on each of which the support is smooth, as for verified basis
regions with a fixed or a rebuilt matrix.

\begin{theorem}[Comparison matrix from a complete region cover]\label{thm:cover-con}
Let $\mathcal D\subseteq\R^k$ be closed and convex, and let
$F:\mathcal D\to\R^{2n}$. For each $i\in\{1,\ldots,2n\}$ let
$C_{i,1},\ldots,C_{i,N_i}$ be closed sets whose union contains $\mathcal D$,
such that every line segment contained in $\mathcal D$ meets each
$C_{i,s}$ in a finite union of closed intervals, single points allowed.
Suppose that for each $(i,s)$ there is a function $h_{i,s}$, continuously
differentiable on an open set containing $C_{i,s}\cap\mathcal D$, with
$F_i=h_{i,s}$ on $C_{i,s}\cap\mathcal D$. If $M\ge0$ satisfies
\[
 |\partial_jh_{i,s}(\theta)|\le M_{ij}
 \qquad(\theta\in C_{i,s}\cap\mathcal D,\ \text{all }i,s,j),
\]
then $|F(p)-F(q)|\le M|p-q|$ for all $p,q\in\mathcal D$. In particular,
$F$ satisfies~\eqref{eq:matrix-lipschitz} on $\mathcal D$.
\end{theorem}

The proof is in Appendix~\ref{sec:parametric-proofs-cover}.

No continuity of $F$ is assumed separately: on each piece of the partition
$F_i$ coincides with one smooth function, and consecutive pieces share their
endpoint values because they evaluate the same function $F_i$. For the
Jacobi map~\eqref{eq:endpoint-map-con}, take $\theta=p$ and, for each
direction $i$, regions on which a basis is verified. With a fixed matrix
these are the polyhedra $C_I$ of Proposition~\ref{prop:basis-region-con},
each meets a segment in one interval, and $M$ must bound $|E^T\pi|$
entrywise for the dual vectors $\pi$ of all bases used. With rebuilt rows,
take closed polyhedra $\Pi$ verified by
Proposition~\ref{prop:rebuilt-basis-con}; then $M$ must
bound~\eqref{eq:lagrangian-derivative-con} over each $\Pi\cap\mathcal D$.
A region may also use any smooth function that agrees with $F_i$ on it;
this handles parameters such as degenerate boxes, at which a basis matrix
becomes singular but the support is known directly.

Coverage of $\mathcal D$ by finitely many polyhedra is a separate
obligation. In one parameter, sorted interval endpoints decide it. In
general it can be decided by finitely many exact linear programs.

\begin{proposition}[A finite coverage test]\label{prop:coverage-con}
Let $\mathcal D\subset\R^k$ be a nonempty polytope, and let
$C^r=\{\theta:H^r\theta\le g^r\}$, $r=1,\ldots,N_C$, be polyhedra, each
given by at least one inequality. For every tuple
$J=(j_1,\ldots,j_{N_C})$ that selects one inequality $j_r$ of each $C^r$,
let
\[
 \tau_J=\max\bigl\{t:\ \theta\in\mathcal D,\ t\le1,\
  t\le H^r_{j_r}\theta-g^r_{j_r}\ (r=1,\ldots,N_C)\bigr\}.
\]
Then $\mathcal D\subseteq\bigcup_rC^r$ if and only if $\tau_J\le0$ for
every tuple $J$.
\end{proposition}

The proof is in Appendix~\ref{sec:parametric-proofs-coverage}.

An exact dual feasible vector of value at most zero for each tuple
certifies coverage, and a feasible point with $t>0$ exhibits an uncovered
parameter. Shared boundaries give $\tau_J=0$ and are handled correctly.
The number of tuples is the product of the numbers of inequalities, so the
test is a finite verification, not an efficient discovery procedure.
Checking only the vertices of $\mathcal D$ is not sufficient: on
$[-1,1]^2$ the regions $\{\theta_1\le-1/2\}$ and $\{\theta_1\ge1/2\}$
contain all four vertices but miss the strip $|\theta_1|<1/2$, and their one
tuple gives $\tau_J=1/2$. The test applies to each support direction
separately; a complete cover for one direction says nothing about another.

\begin{corollary}[Input for the residual theorem]\label{cor:residual-input-con}
Let $P$ be protected with threshold at most $U$, with endpoint vector
$p_P\le p_0$. If, for each direction $i$, finitely many verified regions
cover $[p_P,p_0]$ and $M$ satisfies the derivative bounds of
Theorem~\ref{thm:cover-con}, then~\eqref{eq:matrix-lipschitz} holds on the
invariant interval $[p_P,p_0]$, and Theorem~\ref{thm:residual}
applies with every $e\ge0$ that satisfies~\eqref{eq:supersolution}.
\end{corollary}

\begin{proof}
Lemma~\ref{lem:invariant-interval} shows that $[p_P,p_0]$ is invariant.
It is closed and convex, so Theorem~\ref{thm:cover-con} applies.
\end{proof}

The supersolution inequality~\eqref{eq:supersolution} is then a finite
exact check. If verified regions cover a convex set of pairs $(p,U)$,
Theorem~\ref{thm:cover-con} applied with $\theta=(p,U)$ gives the stronger
estimate $|F_U(p)-F_V(q)|\le M|p-q|+h|U-V|$ on that set. Because the cutoff
row has fixed coefficients, $h_i$ is a bound on the cutoff multipliers of
the bases used for direction $i$. Such an $h$ also serves in
Proposition~\ref{prop:fixed-point-shift}, which needs it only at one fixed
point.

\begin{example}[A basis change along the trajectory]\label{ex:switch-con}
Minimize $x^2$ subject to $x+x^2\le1$ and $0\le x\le4$. The minimizer is
$x^*=0$ with $f^*=0$; take the cutoff $U=0$. On $B=[0,u]$, relax $s=x^2$
by the McCormick rows $s\ge0$, $s\ge2ux-u^2$, $s\le ux$, and keep the
coupling row $x+s\le1$ and the cutoff $s\le0$. The lower endpoint stays at
$0$, since $(0,0)$ is feasible on every box. The cutoff and $s\ge0$ force
$s=0$, so the tangent row gives $x\le u/2$ and the coupling row gives
$x\le1$. Hence the upper support is
\[
 F(u)=\min\{u/2,\ 1\},\qquad 0\le u\le4 .
\]
For $0<u\le2$ the basis of the tangent row $2ux-s\le u^2$ and the cutoff
has primal point $(u/2,0)$ and multipliers $(1/(2u),1/(2u))$; all other rows
hold because $u/2\le1$. Formula~\eqref{eq:lagrangian-derivative-con} gives
$\partial F/\partial u=(1/(2u))(2u-2\cdot u/2)=1/2$, which agrees with direct
differentiation of $u/2$. At $u=0$ the box is $\{0\}$ and $F(0)=0$, so the
smooth function $u/2$ represents $F$ on the closed region $[0,2]$. For
$2\le u\le4$ the basis of $x+s\le1$ and $-s\le0$ has primal point $(1,0)$
and multipliers $(1,1)$; the tangent row holds because $2u\le u^2$. Its
rows do not depend on $u$, so the slope is $0$.

The degenerate box $P=\{0\}$ is protected at cutoff $0$, with the lifted
witness $(0,0)$, and Corollary~\ref{cor:residual-input-con} applies on
$\{(0,u):0\le u\le4\}$, which the two regions cover. Every point of this
interval has $\ell=0$, and the lower support is identically $0$. Hence
$M=\operatorname{diag}(0,1/2)$ satisfies~\eqref{eq:matrix-lipschitz} on it.
From $u_0=4$, the exact iterates are $4,1,1/2,1/4,\ldots$. The first
residual is $d=(0,3)$, and $e=(0,6)$ satisfies~\eqref{eq:supersolution}.
Theorem~\ref{thm:residual} bounds the total motion by $6$ and the
motion after the first round by $Me=3$; the actual values are $4$ and $1$.

The observed displacements $3$ and $1/2$ have ratio $1/6$. A geometric
extrapolation with that ratio would bound the motion after the first round
by $(1/2)/(1-1/6)=3/5$, but the actual motion is $1$: the first round used
the coupling row, and later rounds use the tangent row with ratio $1/2$.
This is a McCormick instance of the failure shown abstractly in
Example~\ref{ex:ratio} and Proposition~\ref{prop:history}. The frozen rows
of $[0,4]$ give the support $1$ for every $u\in[1,4]$, so their comparison
matrix on that interval is zero. Used in Theorem~\ref{thm:residual},
it would assert that no motion remains after the first round, whereas the
rebuilt iterates continue to $0$.
\end{example}

Complete covers are practical when the parameter set is low-dimensional,
for example a scalar cutoff or a few endpoints that actually move. The
number of basis regions is bounded only by the number of bases, which grows
combinatorially with the number of rows, and the coverage test of
Proposition~\ref{prop:coverage-con} has a combinatorial number of tuples.
Theorem~\ref{thm:cover-con} is therefore a verifiable sufficient condition,
not an inexpensive general method.

\subsection{An affine equality changes the rate}\label{sec:equality-con}

Assumption~\ref{ass:tangent} requires finite relaxation values on the
comparison boxes. A projected relaxation that keeps an equality constraint
exactly is infinite off the feasible affine set, so the constrained case
needs a restricted tangent construction, as observed in
Proposition~\ref{prop:restricted-tangent-con}. The following example
computes the restricted map exactly.

Consider $f(x,y)=(x-y)^2=x^2+y^2-2xy$ on boxes $[-r,r]^2$. Keep the squares
exact and replace $xy$ by its McCormick overestimator
$\min\{r^2+r(x-y),\,r^2-r(x-y)\}=r^2-r|x-y|$. The relaxed objective is
\[
 \phi_r(x,y)=x^2+y^2-2r^2+2r|x-y|.
\]
Without further constraints, the points $(r,r)$ and $(-r,-r)$ have
$\phi_r=0$ and are true minimizers, so at cutoff $0$ every box face is
retained and OBBT stalls.

\begin{proposition}[Exact map with an affine equality]\label{prop:equality-con}
Impose the equality $y=-x$ exactly in the relaxation, and let $U\ge0$.
A Jacobi round maps $[-r,r]^2$ to $[-G_U(r),G_U(r)]^2$, where
\[
 G_U(r)=\min\bigl\{r,\ \sqrt{2r^2+U/2}-r\bigr\}.
\]
Put $a=\sqrt U/2$ and $r_{k+1}=G_U(r_k)$.
\begin{enumerate}
\item[(a)] If $U=0$, then $r_{k+1}=(\sqrt2-1)r_k$.
\item[(b)] If $U>0$ and $r_0>a$, then $r_k$ decreases strictly to $a$, and
 \[
  r_{k+1}-a=\frac{(r_k-a)^2}{\sqrt{2r_k^2+2a^2}+r_k+a}
  \le\min\Bigl\{\frac{r_k-a}{2},\ \frac{(r_k-a)^2}{4a}\Bigr\}.
 \]
\item[(c)] If $U>0$ and $r_0\le a$, then $r_k=r_0$ for all $k$.
\end{enumerate}
\end{proposition}

\begin{proof}
On the equality put $(x,y)=(t,-t)$. Then
$\phi_r(t,-t)=2t^2+4r|t|-2r^2$, and the cutoff
$2|t|^2+4r|t|-2r^2\le U$ holds exactly when
$|t|\le\sqrt{2r^2+U/2}-r$, the nonnegative root. The box adds $|t|\le r$.
Both coordinates range over $[-G_U(r),G_U(r)]$, which proves the formula
for the map. For $U=0$ the root is $(\sqrt2-1)r$.

For $U>0$ we have $U/2=2a^2$. Put $X=2r^2+2a^2$. Since
$X-(r+a)^2=(r-a)^2\ge0$, we get $\sqrt X\ge r+a$ and
\[
 \sqrt X-r-a=\frac{X-(r+a)^2}{\sqrt X+r+a}=\frac{(r-a)^2}{\sqrt X+r+a}.
\]
The root is smaller than $r$ exactly when $X<4r^2$, that is, when $r>a$.
Hence for $r>a$ the map equals the root, lies in $(a,r)$, and satisfies
the displayed identity. The denominator is at least $2(r+a)$, which is at
least both $2(r-a)$ and $4a$; this gives the two bounds, and with them
strict decrease to $a$. If $r\le a$, the root is at least $r$, so the box
is fixed.
\end{proof}

For $U>0$ and every starting radius $r_0$, the limiting radius is
$\min\{r_0,a\}$. This is the radius of the original cutoff set within the
initial box: on $y=-x$ the objective is $4x^2$, and $4x^2\le U$ exactly
when $|x|\le a$. OBBT therefore converges to the box hull of the original
cutoff set. When $r_0>a$ it does not reach it in finitely many rounds, since
$r_{k+1}-a>0$ whenever $r_k>a$, but the error decreases quadratically,
whereas in the unconstrained
two-variable example the error at a positive cutoff decreases only
linearly (Proposition~\ref{prop:two-variable-cutoff-floor}(d)). At zero cutoff the
contraction factor is $\sqrt2-1$. Eliminating the equality before relaxing
gives the convex objective $4x^2$, whose exact relaxation reaches radius
$\min\{r,a\}$ in one round. The rate is a property of the formulation and
the relaxation, not of the constrained problem alone.

Proposition~\ref{prop:restricted-tangent-con} in
Section~\ref{sec:local-rates-contraction} gives the general restricted tangent
contraction result. Nonlinear equalities require control of their curvature
and of the relaxation remainder, as in the feasible-repair analysis below.

\subsection{Nonlinear constraints through a feasible repair}\label{sec:repair-con}

For nonlinear constraints, a width bound can be obtained without guessing
which constraints will be active at the next support solutions. Each
relaxed point is moved to a nearby original feasible point, and feasible
quadratic growth then bounds its distance to the minimizer. This transfers
a growth condition on feasible points to infeasible relaxed points through
a constraint error bound, as in exact-penalty arguments. The OBBT content
is the explicit list of uniform constants and their consequence for the
iteration.

Let $S=X\cap B_0$ be the original feasible set, including any
integrality, and let $x^*\in S$ be a global minimizer with value $f^*$. Let
$\mathfrak B^*\subseteq\mathfrak B$ be a set of boxes such that
$x^*\in B$ and $\phi_B(x^*)\le f^*$ for every $B\in\mathfrak B^*$, and fix
$U=f^*+\epsilon$ with $\epsilon\ge0$. Let $\nu\ge0$ be a residual
function, defined on the boxes of $\mathfrak B^*$ and zero on $S$, such as
the maximum violation of the original constraints. Assume constants
$\tau,\beta,\kappa,L\ge0$, $\mu>0$, and $\alpha\in(0,1]$ such that every
$B\in\mathfrak B^*$ and every $x\in K_U(B)$ satisfy
\begin{align}
 \phi_B(x)&\ge f(x)-\tau w(B)^2,\qquad \nu(x)\le\beta w(B)^2,
 \label{eq:repair-error-con}
\end{align}
and $x$ has a \emph{feasible repair} $y\in S$ with
\begin{align}
 \norm{x-y}_\infty&\le\kappa\,\nu(x)^\alpha,\label{eq:repair-distance-con}\\
 f(y)&\le f(x)+L\norm{x-y}_\infty,\label{eq:repair-lipschitz-con}\\
 f(y)&\ge f^*+\mu\norm{y-x^*}_\infty^2.\label{eq:repair-growth-con}
\end{align}

Inequality~\eqref{eq:repair-distance-con} is a H\"older error bound for the
constraint system at $x$. For a polyhedral feasible set
$S=\{x:Gx\le g\}$ and $\nu(x)=\max_j(G_jx-g_j)_+$, Hoffman's
theorem~\cite{hoffman1952-on-approximate-solutions-of-systems} gives such
a bound with $\alpha=1$ and a constant $\kappa$ that depends only on $G$.
Inequality~\eqref{eq:repair-lipschitz-con} holds if $\norm{\nabla f}_1\le L$
on a convex set containing $x$ and $y$; only the one-sided estimate on the
repair pair is needed. Inequality~\eqref{eq:repair-growth-con} is feasible
quadratic growth. The repair need not lie in $B$. If every box of
$\mathfrak B^*$ has width at most $\bar w$, every repair lies within
$\rho=\bar w+\kappa(\beta\bar w^2)^\alpha$ of $x^*$, so it suffices that the
Lipschitz and growth estimates hold on that neighborhood. These are
assumptions about every relaxed point that satisfies the cutoff, including
points that violate the original constraints. A small KKT residual at a
support optimizer does not establish them. Without constraints one may take
$\nu\equiv0$ and $y=x$; then~\eqref{eq:repair-error-con}
and~\eqref{eq:repair-growth-con} combine to the hypothesis
$\phi_B(x)\ge f^*+\mu\norm{x-x^*}_\infty^2-\tau w(B)^2$ of
Proposition~\ref{prop:quadratic-growth}, and the next theorem reduces to it.

\begin{theorem}[Width bound from a feasible repair]\label{thm:repair-con}
For every $B\in\mathfrak B^*$, with $w=w(B)$ and $r=\kappa(\beta w^2)^\alpha$,
\begin{equation}\label{eq:repair-width-con}
 w(T_U(B))\le2r+2\sqrt{\frac{\epsilon+\tau w^2+Lr}{\mu}}.
\end{equation}
\end{theorem}

\begin{proof}
Validity at $x^*$ gives $x^*\in K_U(B)$, so the set is nonempty. Let
$x\in K_U(B)$ with repair $y$. By~\eqref{eq:repair-error-con},
$f(x)\le\phi_B(x)+\tau w^2\le f^*+\epsilon+\tau w^2$ and
$\nu(x)\le\beta w^2$, so~\eqref{eq:repair-distance-con} gives
$\norm{x-y}_\infty\le r$. Then~\eqref{eq:repair-lipschitz-con} gives
$f(y)\le f^*+\epsilon+\tau w^2+Lr$, and~\eqref{eq:repair-growth-con} gives
$\norm{y-x^*}_\infty\le\sqrt{(\epsilon+\tau w^2+Lr)/\mu}$. By the triangle
inequality every point of $K_U(B)$ lies within $r$ plus this radius of
$x^*$ in every coordinate, so each coordinate range of $T_U(B)$ is at most
twice that amount.
\end{proof}

\begin{corollary}[Linear repair]\label{cor:repair-rate-con}
Let $\alpha=1$, put $K=\kappa\beta$ and
$\lambda_0=2\sqrt{(\tau+LK)/\mu}$. For every $B\in\mathfrak B^*$,
\begin{equation}\label{eq:repair-step-con}
 w(T_U(B))\le\bigl(\lambda_0+2Kw(B)\bigr)w(B)+2\sqrt{\epsilon/\mu}.
\end{equation}
Suppose $\mathfrak B^*$ contains every subbox of $B_0$ that contains $x^*$
and has width at most $\bar w$, and $\lambda=\lambda_0+2K\bar w<1$. Then
Jacobi iterates $B_{k+1}=T_U(B_k)$ starting from such a box $B_0'$ satisfy
\begin{equation}\label{eq:repair-iterates-con}
 w(B_k)\le\lambda^kw(B_0')+2\sqrt{\epsilon/\mu}\,\frac{1-\lambda^k}{1-\lambda},
 \qquad
 \limsup_{k\to\infty}w(B_k)\le\frac{2\sqrt{\epsilon/\mu}}{1-\lambda}.
\end{equation}
The same bounds hold for any update that replaces $B_k$ by a box containing
$x^*$ and contained in $T_U(B_k)$, such as a complete sequential round
(Lemma~\ref{lem:sequential}(b)).
\end{corollary}

\begin{proof}
With $r=Kw^2$, inequality~\eqref{eq:repair-width-con} and
$\sqrt{s+t}\le\sqrt s+\sqrt t$ give~\eqref{eq:repair-step-con}. The
iterates are nested, contain $x^*$ by validity, and therefore remain in
$\mathfrak B^*$ with width at most $\bar w$; so
$w(B_{k+1})\le\lambda w(B_k)+2\sqrt{\epsilon/\mu}$. Induction proves the
first bound in~\eqref{eq:repair-iterates-con}, and letting $k\to\infty$
proves the second. The final statement uses only the containment in
$T_U(B_k)$.
\end{proof}

At zero cutoff slack the widths contract by at least the factor $\lambda$
in every round, for every box shape in $\mathfrak B^*$, whichever
constraints are active at the support solutions. At positive slack the
limit bound is an upper bound, not an exact floor. The corollary is a
statement of \emph{guaranteed progress}: whenever
$w(B_k)>2\sqrt{\epsilon/\mu}/(1-\lambda)$, the next round reduces the width
by at least $(1-\lambda)w(B_k)-2\sqrt{\epsilon/\mu}>0$. It is therefore
complementary to the protected boxes of Section~\ref{sec:certificates} and
to Theorem~\ref{thm:residual}, which bound the remaining benefit from
above. The two are consistent: a protected box $P\in\mathfrak B^*$
satisfies $w(P)=w(T_U(P))\le\lambda w(P)+2\sqrt{\epsilon/\mu}$, so its
width is at most $2\sqrt{\epsilon/\mu}/(1-\lambda)$, and at zero slack
$P=\{x^*\}$. The constants are analytic premises. They involve $x^*$ and
$f^*$ and are not computed by any procedure in
Section~\ref{sec:algorithms}.

For $\alpha<1$ write $r=\kappa\beta^\alpha w^{2\alpha}$. At zero slack and
$L\kappa\beta>0$, the bound~\eqref{eq:repair-width-con} contains the term
$2\sqrt{L\kappa\beta^\alpha/\mu}\,w^\alpha$, whose ratio to $w$ is unbounded
as $w\downarrow0$, so this argument does not certify linear contraction. It
does not prove a stall either. A local repair constant $L_B=O(w^{2-2\alpha})$
would make the square-root term $O(w)$, but the term $2r$ is $O(w)$ only if
$\alpha\ge1/2$; for $\alpha<1/2$ the argument needs a stronger repair estimate.

\begin{example}[A nonlinear equality with explicit constants]\label{ex:graph-con}
Let
\[
 S=\{(x,y):y=x^2\}\cap\bigl([-1/4,1/4]\times[0,1/16]\bigr),\qquad
 f(x,y)=x^2+y^2-xy/16 .
\]
Both $x$ and $y$ are original variables. For boxes
$B=[\ell,u]\times[0,\bar y]$ with $-1/4\le\ell\le0\le u\le1/4$ and
$0\le\bar y\le1/16$, use the box rows, the graph relaxation
\[
 x^2\le y\le(\ell+u)x-\ell u,
\]
and the relaxed objective
\[
 \phi_B(x,y)=x^2+y^2-\tfrac1{16}m_B(x,y),\qquad
 m_B(x,y)=\min\{uy,\ \ell y+\bar y(x-\ell)\},
\]
where $m_B$ is the McCormick overestimator of $xy$ on $B$. Then $\phi_B$ is
convex, as the maximum of two convex quadratics, and the relaxed feasible
set is convex.

\emph{Minimizer.} On $S$, $f(x,x^2)=x^2(1-x/16+x^2)\ge(63/64)x^2$, because
$x\le1/4$. Since $|x|\le1/4$ implies $x^2\le|x|$, this equals
$(63/64)\norm{(x,x^2)}_\infty^2$. Hence $(0,0)$ is the unique minimizer,
$f^*=0$, and~\eqref{eq:repair-growth-con} holds on $S$ with $\mu=63/64$.

\emph{Validity and monotonicity.} For $x\in[\ell,u]$,
$(\ell+u)x-\ell u-x^2=(x-\ell)(u-x)\ge0$, so the secant is valid, and
$m_B-xy=\min\{(u-x)y,\ (x-\ell)(\bar y-y)\}\ge0$ on $B$, so $\phi_B\le f$
on $S\cap B$. For a nested box $B'=[\ell',u']\times[0,\bar y']\subseteq B$,
the factors $x-\ell'\le x-\ell$ and $u'-x\le u-x$ show that the new secant
is no higher on $B'$. Also $u'y\le uy$, and
\[
 \ell'y+\bar y'(x-\ell')\le\ell'y+\bar y(x-\ell')\le\ell y+\bar y(x-\ell),
\]
using $x\ge\ell'$ and $\bar y'\le\bar y$ for the first inequality and
$(\ell'-\ell)(y-\bar y)\le0$ for the second. Thus $m_{B'}\le m_B$ and
$\phi_{B'}\ge\phi_B$ on $B'$, and the relaxed feasible set shrinks: the
whole construction is monotone. The origin is relaxed feasible and every
relaxed point has $y\ge0$, so OBBT keeps the lower $y$-bound at $0$, and all
iterates stay in this family of boxes and contain the origin.

\emph{Residual, repair, and objective error.} Let $\nu(x,y)=|y-x^2|$. Every
relaxed point has the repair $(x,x^2)\in S$, at distance
$\delta=y-x^2$ in the infinity norm, and
\[
 0\le\delta\le(x-\ell)(u-x)=\tfrac{(u-\ell)^2}4-\bigl(x-\tfrac{\ell+u}2\bigr)^2
 \le\tfrac14w(B)^2 .
\]
So $\beta=1/4$, $\kappa=1$, and $\alpha=1$. With $s_1=u-x$, $s_2=x-\ell$,
$t_1=y$, and $t_2=\bar y-y$, all nonnegative on $B$, the
arithmetic--geometric mean inequality gives
$m_B-xy=\min\{s_1t_1,s_2t_2\}\le\sqrt{s_1s_2t_1t_2}\le(u-\ell)\bar y/4
\le w(B)^2/4$, so $f-\phi_B\le w(B)^2/64$ and $\tau=1/64$. Finally,
\[
 f(x,x^2)-f(x,y)=\delta\Bigl(\frac x{16}-2x^2-\delta\Bigr)
 =\delta\Bigl(\frac1{2048}-2\bigl(x-\tfrac1{64}\bigr)^2-\delta\Bigr)
 \le\frac{\delta}{2048},
\]
so~\eqref{eq:repair-lipschitz-con} holds with $L=1/2048$.

\emph{Rate.} With $K=1/4$, $\lambda_0^2=4(\tau+LK)/\mu=43/672<(13/48)^2$.
Every admitted box has $w(B)\le1/2$, so~\eqref{eq:repair-step-con} gives,
for all admitted boxes,
\[
 w(T_U(B))\le\Bigl(\lambda_0+\frac{w(B)}2\Bigr)w(B)
  +2\sqrt{\frac{64\epsilon}{63}},
\]
with factor at most $13/48+1/4=25/48$ for every admitted box and at most
$13/48+1/16=1/3$ once $w(B)\le1/8$. At zero cutoff,
Corollary~\ref{cor:repair-rate-con} certifies the contraction factor
$25/48$ per round from any admitted starting box, and the factor $1/3$ once
the width is at most $1/8$.
\end{example}

The example certifies a uniform rate across all box shapes without
choosing an active constraint: at a support solution either the graph
epigraph or the secant may be active, and the repair treats every relaxed
point alike. The constants are sufficient, not exact rates. The constraint
need not be relaxed exactly. If $x^2\le y$ is replaced by its endpoint
tangents $y\ge2\ell x-\ell^2$ and $y\ge2ux-u^2$, the construction remains
monotone, because on a smaller interval $[\ell',u']$ the new tangents
dominate the old ones: $(2\ell'x-\ell'^2)-(2\ell x-\ell^2)
=(\ell'-\ell)(2x-\ell'-\ell)\ge(\ell'-\ell)^2$ for $x\ge\ell'$, and
symmetrically at the upper endpoint. Relaxed points may now lie below the
graph, but still $|y-x^2|\le w(B)^2/4$. For $\delta<0$ the same
identity gives $f(x,x^2)-f(x,y)=|\delta|(2x^2-|\delta|-x/16)\le(9/64)|\delta|$,
so $L=9/64$, $\lambda_0^2=13/63<1/4$, and the contraction factor is at most
$3/4$ for every admitted box. The objective squares, however, are kept
exact. If they are also replaced by endpoint tangents, the generic bound on
the error of each square is $w(B)^2/4$; with $\tau\ge1/4$ the constant
$\lambda_0$ exceeds $1$, so this sufficient estimate certifies no
contraction. That does not show that the resulting OBBT operator fails to
contract. The tangent-map analysis of Section~\ref{sec:local-rates} is the
appropriate tool for such finite linear relaxations.

\begin{table}[tbp]
\centering
\small
\caption{Stored information for constrained support problems and the scope
of its conclusions.}\label{tab:scope-con}
\begin{tabular}{@{}>{\raggedright\arraybackslash}p{0.23\textwidth}%
>{\raggedright\arraybackslash}p{0.22\textwidth}%
>{\raggedright\arraybackslash}p{0.24\textwidth}%
>{\raggedright\arraybackslash}p{0.23\textwidth}@{}}
\toprule
Information & Remains valid after & Conclusion & Additional verification\\
\midrule
Dual vector, fixed matrix (Prop.~\ref{prop:dual-envelope-con}) &
 any change of the right-hand-side parameter &
 upper bound on the support; guaranteed reduction &
 exact dual feasibility\\
Optimal dual with cutoff multiplier $\eta$ (Cor.~\ref{cor:cutoff-multiplier-con}) &
 a lower cutoff, same box &
 reduction at least $\eta(U_0-U)$ &
 exact optimality at $U_0$\\
Basis region (Prop.~\ref{prop:basis-region-con}) &
 parameters in the region &
 exact affine support &
 all primal residuals on the region\\
Retained witness (Prop.~\ref{prop:bracket-con}) &
 a cutoff at least its objective, same box and rows &
 lower bound; limits the possible reduction &
 recheck all rows after any other change\\
Least-objective optimal witness (Prop.~\ref{prop:threshold-con}) &
 a lower cutoff, same box &
 exact no-change threshold $\underline U$ &
 one additional solve\\
Frozen-row dual (Lem.~\ref{lem:frozen-con}) &
 a smaller box with rebuilt rows &
 valid bound, but no progress after one round &
 lifted monotonicity\\
Corrected stored dual (Prop.~\ref{prop:corrected-dual-con}) &
 a smaller box with rebuilt rows &
 valid upper bound &
 bounds on columns with changed coefficients\\
Complete region cover (Thm.~\ref{thm:cover-con}) &
 all parameters in $\mathcal D$ &
 comparison matrix for Theorem~\ref{thm:residual} &
 coverage test, derivative bounds, invariance\\
Repair constants (Thm.~\ref{thm:repair-con}) &
 all boxes of $\mathfrak B^*$ &
 width recurrence; guaranteed progress &
 uniform error, repair, and growth bounds\\
\bottomrule
\end{tabular}
\end{table}

\subsection{Summary}\label{sec:summary-con}

Table~\ref{tab:scope-con} lists the information discussed in this section,
the changes after which it remains valid, and what it proves. Observed
progress is not on the list, for the reasons given in
Section~\ref{sec:histories-con}.
